\chapter*{Resumen}
\addcontentsline{toc}{chapter}{Resumen}

\textbf{Problema.} La mayor\'ia de cl\'inicas veterinarias en Ecuador
operan con registros manuales o en papel: p\'erdida de expedientes,
dif\'icil consulta del historial cl\'inico y falta de indicadores de
gesti\'on~\cite{cedeno2021veterinaria}. \textbf{Soluci\'on.} BIOPET es un
sistema web de gesti\'on veterinaria (Spring~Boot, Angular, PostgreSQL,
Redis/Valkey) con autenticaci\'on por rol/propiedad, CRUD de mascotas,
citas, consultas, vacunas y acceso a rutinas almacenadas.
\textbf{Metodolog\'ia.} \emph{Engineering Research}/\emph{Design
Science} (Peffers et al.~\cite{peffers2007dsrm}) con marco
Goal--Question--Metric~\cite{vansolingen2002gqm}, sin dise\~no
experimental controlado. \textbf{Principales validaciones.} Corte
hist\'orico v1.0.0: 205 pruebas automatizadas (0 fallos, 0 errores) y
cobertura JaCoCo de 91{,}80\,\% en l\'ineas y 79{,}39\,\% en ramas
(umbral 70\,\%); rendimiento con k6
(0{,}0\,\% de error); usabilidad SUS sobre 18 participantes (media
74{,}44/100, IC~95\,\% $[63{,}33, 85{,}56]$); seguridad (OWASP
Top~10 manual; SpotBugs est\'atico y auditor\'ia de SQL din\'amico sin
hallazgos; ZAP: \emph{baseline} hist\'orico y escaneo autenticado sin
alertas altas); y despliegue en producci\'on (Render,
HTTPS activo). Software, evidencias y backend archivados con DOI
(\texttt{\doisoftware}/\texttt{\doidataset}, Zenodo) y \emph{digest}
inmutable (GHCR). \textbf{Conclusi\'on
principal.} BIOPET demuestra, con evidencia reproducible
(ISO/IEC~25010~\cite{iso25010-2011}), que es posible construir y evaluar
emp\'iricamente, en un Proyecto Fin de Curso,
un artefacto funcional, seguro seg\'un los controles auditados y
reproducible; las limitaciones (entorno acad\'emico, tama\~no muestral)
se documentan en el cap\'itulo de amenazas a la validez.

\textbf{Palabras clave:} gesti\'on veterinaria; aplicaciones web;
Spring~Boot; Angular; JWT; PostgreSQL; ISO/IEC~25010; usabilidad (SUS);
seguridad OWASP; reproducibilidad.

\bigskip
\hrule
\bigskip

\chapter*{Abstract}
\addcontentsline{toc}{chapter}{Abstract}

\textbf{Problem.} Most veterinary clinics in Ecuador and similar contexts
still rely on manual or paper-based records: lost files, difficult
clinical-history retrieval, and a lack of management
indicators~\cite{cedeno2021veterinaria}. \textbf{Solution.} BIOPET is a
web-based veterinary management system (Spring~Boot, Angular,
PostgreSQL, Redis/Valkey) with role/ownership-based authentication, pet
CRUD, appointments, consultations, vaccines, and access to stored
routines. \textbf{Methodology.} Engineering Research/Design Science
(Peffers et al.~\cite{peffers2007dsrm}) with a Goal--Question--Metric
framework~\cite{vansolingen2002gqm}, no controlled experimental design.
\textbf{Main validations.} In the historical v1.0.0 snapshot, 205
automated tests were recorded (0 failures, 0 errors), with JaCoCo line
coverage of 91.80\,\% and branch coverage of 79.39\,\% (70\,\% gate); k6
performance testing (0.0\,\% error rate); usability measured
with SUS on 18 participants (mean 74.44/100, 95\,\% CI $[63.33, 85.56]$);
security audited (manual OWASP Top~10; SpotBugs static analysis and a
dynamic-SQL audit of the stored routines, both with no findings; dynamic
OWASP~ZAP: a historical baseline scan and an authenticated scan on the
local TLS environment, both with no high-severity alerts); and
production deployment (Render, HTTPS enabled). Software, evidence and
backend image remain archived with a DOI
(\texttt{\doisoftware}/\texttt{\doidataset}, Zenodo) and an immutable
\emph{digest} (GHCR). \textbf{Main conclusion.}
BIOPET shows, with reproducible evidence (ISO/IEC~25010~\cite{iso25010-2011}),
that it is possible to build and empirically evaluate, within a capstone
course project, a functional artifact that is secure with respect to the
audited controls and reproducible; limitations (academic environment,
sample size) are documented in the threats-to-validity chapter.

\textbf{Keywords:} veterinary management; web applications;
Spring~Boot; Angular; JWT; PostgreSQL; ISO/IEC~25010; usability (SUS);
OWASP security; reproducibility.
